% GENERATED FILE. Edit canonical lessons, then run npm run generate:books.
% canonical-source-hash: fnv1a64:00d0f38347a3f8de
% canonical-lessons: GU-C521-melvan, GU-C521-jamavvu, GU-C521-valovvu, GU-C521-ubhravu, GU-C521-phinvu

\chapter{Milk, Curd and Butter}
\label{ch:gu-milk-curd-and-butter}
\hlchaptermodality{521}

\begin{chapteropening}
\textbf{By the end of this chapter:} \emph{I can talk about setting curd and churning butter.}

Complete the last lesson of chapter 521: I can talk about setting curd and churning butter.
\end{chapteropening}

\section[\texorpdfstring{\gu{મેળવણ} (meḷvaṇ)}{મેળવણ (meḷvaṇ)}]{\gu{મેળવણ} (meḷvaṇ) — a curd starter}

\label{lesson:GU-C521-melvan}



\begin{quote}
Before the new one: say the Gujarati for to sprinkle, then the Gujarati for to parch.
\end{quote}

\subsection*{You'll want to know: \gu{મેળવણ}}
\textbf{\gu{મેળવણ}} — \emph{meḷvaṇ} — \textquotedblleft{}a curd starter\textquotedblright{}.

\begin{quote}
\gu{દહીં} \gu{માટે} \gu{દૂધમાં} \gu{એક} \gu{ચમચી} \gu{મેળવણ} \gu{નાખો}.
\end{quote}

\begin{quote}
\emph{dahī̃ māṭe dūdhmā̃ ek camcī meḷvaṇ nākho.}
\end{quote}

\begin{quote}
For curd, put a spoonful of starter in the milk.
\end{quote}

A neuter noun.

\begin{cousinweb}
\textbf{\gu{મેળવણ}} (\emph{meḷvaṇ}) is built on \textbf{\gu{મેળવવું}}, \textquotedblleft{}to get; to mix in\textquotedblright{}, with the ending \emph{-aṇ}: the spoonful of old curd mixed into milk to set it. CDIAL files the family under Sanskrit \emph{miláti}, \textquotedblleft{}meets\textquotedblright{} (CDIAL 10133).
\end{cousinweb}

\begin{grammarlens}[title={What you've built}]
One more word for describing what you lived and felt.
\end{grammarlens}

\subsection*{Your turn}
\begin{itemize}
\raggedright
  \item \emph{Say it:} \emph{meḷvaṇ}
  \item \emph{Say it:} \emph{meḷvaṇ}, once more
  \item \emph{Say it:} \emph{dahī̃ māṭe dūdhmā̃ ek camcī meḷvaṇ nākho.} — For curd, put a spoonful of starter in the milk.
  \item \emph{Recall:} say the Gujarati for to sprinkle, then the Gujarati for to parch, then say \emph{meḷvaṇ} again
\end{itemize}

\begin{tcolorbox}[breakable,colback=teal!4,colframe=teal!35!black,title={Before you move on}]
What does \gu{મેળવણ} mean? (\textquotedblleft{}A curd starter\textquotedblright{}.) Say it once more.
\end{tcolorbox}
\section[\texorpdfstring{\gu{જમાવવું} (jamāvvũ)}{જમાવવું (jamāvvũ)}]{\gu{જમાવવું} (jamāvvũ) — to set curd}

\label{lesson:GU-C521-jamavvu}



\begin{quote}
Before the new one: say the Gujarati for to parch, then the Gujarati for a curd starter.
\end{quote}

\subsection*{You'll want to know: \gu{જમાવવું}}
\textbf{\gu{જમાવવું}} — \emph{jamāvvũ} — \textquotedblleft{}to set curd\textquotedblright{}.

\begin{quote}
\gu{દાદીએ} \gu{રાત્રે} \gu{દહીં} \gu{જમાવ્યું}.
\end{quote}

\begin{quote}
\emph{dādīe rātre dahī̃ jamāvyũ.}
\end{quote}

\begin{quote}
Grandmother set the curd at night.
\end{quote}

Past, with \gu{તેણે} (he or she …): \gu{તેણે} \gu{દહીં} \gu{જમાવ્યું}, \emph{teṇe dahī̃ jamāvyũ}.

\begin{cousinweb}
\textbf{\gu{જમાવવું}} (\emph{jamāvvũ}) is filed by CDIAL under Sanskrit \emph{yamyátē}, \textquotedblleft{}is fixed\textquotedblright{}, and CDIAL lists the Gujarati verb as \textquotedblleft{}to collect\textquotedblright{} (CDIAL 10428); in the kitchen it is what you do to milk to make curd.
\end{cousinweb}

\begin{grammarlens}[title={What you've built}]
One more word for describing what you lived and felt.
\end{grammarlens}

\subsection*{Your turn}
\begin{itemize}
\raggedright
  \item \emph{Say it:} \emph{jamāvvũ}
  \item \emph{Say it:} \emph{jamāvvũ}, once more
  \item \emph{Say it:} \emph{dādīe rātre dahī̃ jamāvyũ.} — Grandmother set the curd at night.
  \item \emph{Recall:} say the Gujarati for to parch, then the Gujarati for a curd starter, then say \emph{jamāvvũ} again
\end{itemize}

\begin{tcolorbox}[breakable,colback=teal!4,colframe=teal!35!black,title={Before you move on}]
What does \gu{જમાવવું} mean? (\textquotedblleft{}To set curd\textquotedblright{}.) Say it once more.
\end{tcolorbox}
\section[\texorpdfstring{\gu{વલોવવું} (valovvũ)}{વલોવવું (valovvũ)}]{\gu{વલોવવું} (valovvũ) — to churn}

\label{lesson:GU-C521-valovvu}



\begin{quote}
Before the new one: say the Gujarati for a curd starter, then the Gujarati for to set curd.
\end{quote}

\subsection*{You'll want to know: \gu{વલોવવું}}
\textbf{\gu{વલોવવું}} — \emph{valovvũ} — \textquotedblleft{}to churn\textquotedblright{}.

\begin{quote}
\gu{સવારે} \gu{દાદીએ} \gu{દહીં} \gu{વલોવીને} \gu{માખણ} \gu{કાઢ્યું}.
\end{quote}

\begin{quote}
\emph{savāre dādīe dahī̃ valovīne mākhaṇ kāḍhyũ.}
\end{quote}

\begin{quote}
In the morning my grandmother churned the curd and took out the butter.
\end{quote}

Past, with \gu{તેણે} (he or she …): \gu{તેણે} \gu{દહીં} \gu{વલોવ્યું}, \emph{teṇe dahī̃ valovyũ}.

\begin{cousinweb}
\textbf{\gu{વલોવવું}} (\emph{valovvũ}) is filed by CDIAL under Sanskrit \emph{vilōbhayati}, \textquotedblleft{}perplexes, agitates\textquotedblright{}, and CDIAL lists the Gujarati verb, \textquotedblleft{}to churn\textquotedblright{} (CDIAL 11913a).
\end{cousinweb}

\begin{grammarlens}[title={What you've built}]
One more word for describing what you lived and felt.
\end{grammarlens}

\subsection*{Your turn}
\begin{itemize}
\raggedright
  \item \emph{Say it:} \emph{valovvũ}
  \item \emph{Say it:} \emph{valovvũ}, once more
  \item \emph{Say it:} \emph{savāre dādīe dahī̃ valovīne mākhaṇ kāḍhyũ.} — In the morning my grandmother churned the curd and took out the butter.
  \item \emph{Recall:} say the Gujarati for a curd starter, then the Gujarati for to set curd, then say \emph{valovvũ} again
\end{itemize}

\begin{tcolorbox}[breakable,colback=teal!4,colframe=teal!35!black,title={Before you move on}]
What does \gu{વલોવવું} mean? (\textquotedblleft{}To churn\textquotedblright{}.) Say it once more.
\end{tcolorbox}
\section[\texorpdfstring{\gu{ઊભરાવું} (ūbhrāvũ)}{ઊભરાવું (ūbhrāvũ)}]{\gu{ઊભરાવું} (ūbhrāvũ) — to boil over}

\label{lesson:GU-C521-ubhravu}



\begin{quote}
Before the new one: say the Gujarati for to set curd, then the Gujarati for to churn.
\end{quote}

\subsection*{You'll want to know: \gu{ઊભરાવું}}
\textbf{\gu{ઊભરાવું}} — \emph{ūbhrāvũ} — \textquotedblleft{}to boil over\textquotedblright{}.

\begin{quote}
\gu{દૂધ} \gu{ઊભરાઈ} \gu{ગયું} \gu{અને} \gu{ચૂલો} \gu{ગંદો} \gu{થઈ} \gu{ગયો}.
\end{quote}

\begin{quote}
\emph{dūdh ūbhrāī gayũ ane cūlo gando thaī gayo.}
\end{quote}

\begin{quote}
The milk boiled over and the stove got dirty.
\end{quote}

Past (it …): \gu{તે} \gu{દૂધ} \gu{ઊભરાઈ} \gu{ગયું}, \emph{te dūdh ūbhrāī gayũ}.

\begin{cousinweb}
\textbf{\gu{ઊભરાવું}} (\emph{ūbhrāvũ}) is filed by CDIAL under Sanskrit \emph{údbharati}, \textquotedblleft{}raises\textquotedblright{}, with the Gujarati verb listed as \emph{ubhrāvũ}, \textquotedblleft{}to effervesce\textquotedblright{} (CDIAL 2038).
\end{cousinweb}

\begin{grammarlens}[title={What you've built}]
One more word for describing what you lived and felt.
\end{grammarlens}

\subsection*{Your turn}
\begin{itemize}
\raggedright
  \item \emph{Say it:} \emph{ūbhrāvũ}
  \item \emph{Say it:} \emph{ūbhrāvũ}, once more
  \item \emph{Say it:} \emph{dūdh ūbhrāī gayũ ane cūlo gando thaī gayo.} — The milk boiled over and the stove got dirty.
  \item \emph{Recall:} say the Gujarati for to set curd, then the Gujarati for to churn, then say \emph{ūbhrāvũ} again
\end{itemize}

\begin{tcolorbox}[breakable,colback=teal!4,colframe=teal!35!black,title={Before you move on}]
What does \gu{ઊભરાવું} mean? (\textquotedblleft{}To boil over\textquotedblright{}.) Say it once more.
\end{tcolorbox}
\section[\texorpdfstring{\gu{ફીણવું} (phīṇvũ)}{ફીણવું (phīṇvũ)}]{\gu{ફીણવું} (phīṇvũ) — to whisk (to a froth)}

\label{lesson:GU-C521-phinvu}



\begin{quote}
Before the new one: say the Gujarati for to churn, then the Gujarati for to boil over.
\end{quote}

\subsection*{You'll want to know: \gu{ફીણવું}}
\textbf{\gu{ફીણવું}} — \emph{phīṇvũ} — \textquotedblleft{}to whisk (to a froth)\textquotedblright{}.

\begin{quote}
\gu{મીરાએ} \gu{ચમચીથી} \gu{છાશ} \gu{ફીણી}.
\end{quote}

\begin{quote}
\emph{Mīrāe camcīthī chāś phīṇī.}
\end{quote}

\begin{quote}
Mira whisked the buttermilk with a spoon.
\end{quote}

Past, with \gu{તેણે} (he or she …): \gu{તેણે} \gu{છાશ} \gu{ફીણી}, \emph{teṇe chāś phīṇī}.

\begin{cousinweb}
\textbf{\gu{ફીણવું}} (\emph{phīṇvũ}) is built on \emph{phīṇ}, \textquotedblleft{}foam\textquotedblright{}, which CDIAL lists under Sanskrit \emph{ph\'{\={e}}na}, \textquotedblleft{}foam\textquotedblright{} (CDIAL 9108): to beat a liquid until it froths.
\end{cousinweb}

\begin{grammarlens}[title={What you've built}]
One more word for describing what you lived and felt.
\end{grammarlens}

\subsection*{Your turn}
\begin{itemize}
\raggedright
  \item \emph{Say it:} \emph{phīṇvũ}
  \item \emph{Say it:} \emph{phīṇvũ}, once more
  \item \emph{Say it:} \emph{Mīrāe camcīthī chāś phīṇī.} — Mira whisked the buttermilk with a spoon.
  \item \emph{Recall:} say the Gujarati for to churn, then the Gujarati for to boil over, then say \emph{phīṇvũ} again
\end{itemize}

\begin{tcolorbox}[breakable,colback=teal!4,colframe=teal!35!black,title={Before you move on}]
What does \gu{ફીણવું} mean? (\textquotedblleft{}To whisk (to a froth)\textquotedblright{}.) Say it once more.
\end{tcolorbox}
